\small{

\begin{center}
\MakeUppercase{\textbf{Diffusion distance and diffusion maps: theory and applications}}\\
Bachelor's thesis \\
Riki-Taavi Nurm\\
\end{center}

\normalsize{\textbf{Abstract}}\\
In computer vision and data analysis, we are often interested in studying the intrinsic geometry of a dataset independently of the ambient space in which it lives. After transforming the dataset into a Markov graph, we can define the diffusion distance between any two vertices based on how random walks diffuse between them. Using the eigenvalues and eigenvectors of the corresponding Markov transition matrix, we can define diffusion maps which embed the graph into a Euclidean space. The main aim of this thesis is to prove a theorem which states that the diffusion distance on the graph is equivalent, up to a scaling constant, to the Euclidean distance in the embedded space. We illustrate the theorem on synthetic and established datasets. \\
\textbf{CERCS research specialisation:} P170 Computer science, numerical analysis, systems, control.\\
\textbf{Keywords:} diffusion maps, diffusion distance, spectral theorem, Markov chains, dimensionality reduction.\\
}
\vspace{0.5cm}
\small{

\begin{center}
\MakeUppercase{\textbf{difusioonkaugus ja difusioonikaardid: teooria ja rakendused}}\\
Bakalaureusetöö \\
Riki-Taavi Nurm\\
\end{center}

\normalsize{\textbf{Lühikokkuvõte}}\\
Arvutinägemises ja andmeanalüüsis tahame me tihti uurida andmete sisemist geomeetriat sõltumata ruumist, kus nad asuvad. Pärast andmete Markovi graafiks teisendamist on võimalik iga kahe graafi tipu vahel defineerida difusioonkaugus, kasutades juhuslikku ekslemist nende tippude vahel. Markovi üleminekumaatriksi omaväärtuste ja omavektorite abil saame defineerida difusioonikaardid, mis kujutavad graafi eukleidilisse ruumi. Selle töö peaeesmärk on tõestada teoreem, mis väidab, et difusioonkaugus graafil on konstandiga korrutamise täpsuseni võrdne eukleidilise kaugusega kujutusruumis. Me illustreerime seda teoreemi kasutades nii sünteetilisi kui ka juba tuntud andmehulki.\\
\textbf{CERCS teaduseriala:} P170 Arvutiteadus, arvutusmeetodid, süsteemid, juhtimine
(automaatjuhtimisteooria).\\
\textbf{Märksõnad:} difusioonikaardid, difusioonkaugus, spektraalteoreem, Markovi ahelad, mõõtmete vähendamine.\\
}